\title{Chain-of-Thought Prompting Elicits Reasoning in Large Language Models}

\begin{abstract}
We investigate how producing a \textit{chain of thought}—a sequence of intermediate reasoning steps—can greatly enhance the capability of large language models to tackle complex reasoning tasks. Specifically, we demonstrate that such reasoning skills arise naturally in sufficiently large language models through a straightforward technique called \textit{chain-of-thought prompting}, which involves providing a few examples of chains of thought as prompts.

Tests conducted on three large language models indicate that chain-of-thought prompting boosts performance across various arithmetic, commonsense, and symbolic reasoning challenges. The improvements can be substantial. For example, by prompting a PaLM 540B with only eight chain-of-thought exemplars, we achieve state-of-the-art results on the GSM8K benchmark for math word problems, outperforming even a finetuned GPT-3 equipped with a verifier.
\end{abstract}

\section{Introduction}
The field of NLP has recently undergone a transformation driven by language models \citep[][\textit{inter alia}]{peters-etal-2018-deep,devlin-etal-2019-bert,brown2020language}. Increasing the scale of language models has been shown to yield various advantages, including enhanced performance and better sample efficiency \citep[\textit{inter alia}]{kaplan2020scaling,brown2020language}. Nonetheless, simply enlarging model size has not been enough to achieve strong results on difficult tasks like arithmetic, commonsense reasoning, and symbolic reasoning \citep{rae2021scaling}.

This paper examines how the reasoning capabilities of large language models can be activated through a simple approach inspired by two concepts. First, arithmetic reasoning methods can improve by generating natural language explanations that lead to the final answer. Previous research has enabled models to produce natural language intermediate steps either by training from scratch \citep{ling-etal-2017-program} or finetuning pretrained models \citep{cobbe2021training}, alongside neuro-symbolic strategies employing formal languages instead of natural language \citep{roy-roth-2015-solving, chiang-chen-2019-semantically,amini-etal-2019-mathqa, chen2019neural}. Second, large language models provide an exciting opportunity for in-context few-shot learning through \emph{prompting}. Rather than finetuning separate model checkpoints for each new task, one can simply supply the model with a small number of input-output examples illustrating the task. Notably, this method has proven effective for a variety of straightforward question-answering tasks \citep{brown2020language}.

Both of the aforementioned approaches, however, have significant drawbacks. For rationale-augmented training and finetuning techniques, generating a substantial collection of high-quality rationales is expensive and considerably more complex than the straightforward input--output pairs typically used in conventional machine learning. Meanwhile, the conventional few-shot prompting approach employed by \citet{brown2020language} tends to perform inadequately on tasks that demand reasoning skills, and its performance often does not improve significantly even as the scale of the language model increases \citep{rae2021scaling}. In this study, we integrate the advantages of these two methods while circumventing their respective limitations. Specifically, we investigate the capacity of language models to execute few-shot prompting for reasoning tasks by providing prompts composed of triples: $\langle$input, \emph{chain of thought}, output$\rangle$.

A \emph{chain of thought} refers to a sequence of intermediate natural language reasoning steps that culminate in the final answer, and we label this technique as \textit{chain-of-thought prompting}. An example of such a prompt is depicted in \cref{fig:pull-figure}.

We conduct empirical assessments on benchmarks involving arithmetic, commonsense, and symbolic reasoning, demonstrating that chain-of-thought prompting surpasses standard prompting, occasionally by a substantial margin. \cref{fig:pull-bar-chart} showcases one such outcome—on the GSM8K benchmark of math word problems \citep{cobbe2021training}, chain-of-thought prompting using \palm{} 540B significantly outperforms standard prompting and sets a new state-of-the-art record. Employing a prompting-only strategy is crucial since it does not necessitate a large training dataset and allows a single model checkpoint to handle multiple tasks without sacrificing generality. This research highlights how large language models can learn from just a few examples using natural language information about the task (in contrast to learning patterns relating inputs and outputs from a large training dataset).

\section{Chain-of-Thought Prompting} When tackling a complex reasoning problem, such as a multi-step math word question, people often reflect on their own thinking process. Typically, the problem is broken down into smaller intermediate steps, each solved sequentially before arriving at the final solution: \textit{``After Jane gives 2 flowers to her mom, she has 10 $\ldots$ then after she gives 3 to her dad, she will have 7 $\ldots$ so the answer is 7.''} The aim of this work is to equip language models with the capacity to produce a similar \textit{chain of thought}—a logical sequence of intermediate reasoning steps culminating in the problem's final answer. We demonstrate that large enough language models can generate these chains of thought when examples illustrating chain-of-thought reasoning are included in few-shot prompt demonstrations.

An illustration in \cref{fig:pull-figure} depicts a model generating a chain of thought to resolve a math word problem that it would otherwise answer incorrectly. This chain of thought functions like a solution and can be interpreted as such, but we prefer to term it a chain of thought to emphasize that it emulates a step-by-step reasoning process leading to the answer (additionally, solutions or explanations generally appear \textit{after} the final answer \citep[][\textit{inter alia}]{narang2020wt5,wiegreffe2021reframing,lampinen2022can}).

Chain-of-thought prompting offers several appealing advantages as a method for enhancing reasoning abilities in language models.

\begin{enumerate}[topsep=1pt,itemsep=0ex]
    \item To begin with, chain of thought enables models to break down complex multi-step problems into smaller intermediate steps, allowing additional computational effort to be focused on problems that demand more reasoning stages.
    \item Next, a chain of thought offers a transparent insight into the model’s behavior, indicating the possible reasoning process behind a given answer and enabling debugging of errors in the reasoning path (although fully understanding the computations behind a model’s answer remains an open challenge).
    \item Additionally, chain-of-thought reasoning can be applied to tasks such as mathematical word problems, commonsense reasoning, and symbolic manipulation, and is theoretically extendable to any task solvable by humans through language.
    \item Lastly, chain-of-thought reasoning can be effectively prompted in sufficiently large pre-trained language models simply by incorporating examples of chain of thought sequences into the few-shot prompt exemplars.
\end{enumerate}

In our empirical studies, we will demonstrate the benefits of chain-of-thought prompting across arithmetic reasoning (\cref{sec:arithmetic-reasoning}), commonsense reasoning (\cref{sec:commonsense-reasoning}), and symbolic reasoning (\cref{sec:symbolic-reasoning}).

\section{Arithmetic Reasoning}\label{sec:arithmetic-reasoning}
We start by examining math word problems of the type illustrated in \cref{fig:pull-figure}, which evaluate the arithmetic reasoning capabilities of language models. Although straightforward for humans, arithmetic reasoning often poses difficulties for language models \citep[][\textit{inter alia}]{hendrycks2021measuring,patel-etal-2021-nlp}. Remarkably, chain-of-thought prompting paired with the 540B parameter language model performs on par with task-specific fine-tuned models across several benchmarks, even setting a new state of the art on the challenging GSM8K dataset \citep{cobbe2021training}.

\subsection{Experimental Setup}

We investigate chain-of-thought prompting on multiple language models across various benchmarks.

\textbf{Benchmarks.}
We evaluate on five math word problem datasets:
\textbf{(1)} the \textbf{GSM8K} math word problems benchmark \citep{cobbe2021training},
\textbf{(2)} the \textbf{SVAMP} dataset featuring math word problems with diverse structures \citep{patel-etal-2021-nlp},
\textbf{(3)} the \textbf{ASDiv} collection of varied math word problems \citep{miao-etal-2020-diverse},
\textbf{(4)} the \textbf{AQuA} algebraic word problems dataset, and
\textbf{(5)} the \textbf{MAWPS} benchmark \citep{koncel-kedziorski-etal-2016-mawps}. Sample problems appear in Appendix \cref{tab:math-datasets}.

\textbf{Standard prompting.} As a baseline, we use standard few-shot prompting, popularized by \citet{brown2020language}, where a language model is provided with a set of input--output examples in the context before generating an answer for a test instance. The exemplars are structured as questions followed directly by answers. The model outputs the final answer without intermediate reasoning steps, as illustrated in \cref{fig:pull-figure} (left).

\textbf{Chain-of-thought prompting.} Our method involves enhancing each few-shot exemplar with a chain of thought linked to its corresponding answer, as depicted in \cref{fig:pull-figure} (right). Since most datasets only provide an evaluation split, we manually crafted a set of eight few-shot exemplars that include chains of thought for prompting—\cref{fig:pull-figure} (right) displays one such exemplar, and the complete set is available in Appendix \cref{tab:appendix-mwp-prompts}. (These exemplars were not subjected to prompt engineering; their robustness is examined in \cref{subsec:robustness} and \cref{subsec:faq-prompt-engineering}.) To test whether this style of chain-of-thought prompting can reliably induce effective reasoning over diverse math word problems, we applied the same set of eight exemplars with chains of thought across all benchmarks, except for AQuA, which is a multiple-choice task rather than free response. For AQuA, we employed four exemplars and solutions drawn from the training set, as detailed in Appendix \cref{tab:appendix-aqua-prompt}.

\textbf{Language models.} We assess five large language models. The first is \textbf{GPT-3} \citep{brown2020language}, using variants text-ada-001, text-babbage-001, text-curie-001, and text-davinci-002, which are believed to correspond to InstructGPT models of 350M, 1.3B, 6.7B, and 175B parameters, respectively \citep{ouyang2022training}. The second is \textbf{\lamda{}} \citep{thoppilan2022lamda}, with sizes of 422M, 2B, 8B, 68B, and 137B parameters. The third is \textbf{\palm{}}, comprising models with 8B, 62B, and 540B parameters. The fourth model is \textbf{UL2 20B} \citep{tay2022unifying}, and the fifth is \textbf{Codex} \citep[][code-davinci-002 in the OpenAI API]{chen2021evaluating}. We generate outputs via greedy decoding (although subsequent work indicates that chain-of-thought prompting can be enhanced by taking the majority final answer across multiple sampled generations \citep{wang2022self}). For \lamda{}, results are averaged over five random seeds, each corresponding to a different random shuffle of the exemplar order. Since \lamda{} showed little variance between different seeds, to conserve computational resources, we report results using a single fixed exemplar order for all other models.

\subsection{Results} 
The most notable outcomes of chain-of-thought prompting are presented in \cref{fig:main-math}, with detailed experimental results for each model group, model size, and benchmark provided in \cref{tab:all-lm-math} in the Appendix. There are three primary observations. First, \cref{fig:main-math} illustrates that chain-of-thought prompting emerges as a capability with increasing model scale \citep{wei2022emergent}. Specifically, chain-of-thought prompting does not enhance performance for smaller models and only starts to improve results when applied to models around $\sim$100B parameters. Qualitative analysis showed that smaller models generated coherent but logically flawed chains of thought, resulting in worse performance compared to standard prompting.

\input{main_fables/main-math} 
Second, chain-of-thought prompting yields greater performance improvements on more complex problems. For example, on GSM8K (the dataset with the lowest baseline accuracy), the largest GPT and \palm{} models saw their performance more than double. Conversely, for SingleOp, the simplest subset of MAWPS which requires only one step for solution, improvements were minimal or negative (refer to Appendix \cref{tab:mawps-subsets}).

Third, when applied with GPT-3 175B and \palm{} 540B, chain-of-thought prompting achieves competitive results relative to previous state-of-the-art methods, which usually involve finetuning a model specifically for the task on labeled datasets. \cref{fig:main-math} demonstrates that \palm{} 540B employing chain-of-thought prompting attains new state-of-the-art performance on GSM8K, SVAMP, and MAWPS (noting that standard prompting already outperformed prior best results on SVAMP). For the remaining datasets, AQuA and ASDiv, \palm{} with chain-of-thought prompting approaches the state of the art within 2\% (see Appendix \cref{tab:all-lm-math}).

To gain a deeper understanding of why chain-of-thought prompting is effective, we conducted a manual review of the chains of thought generated by \lamda{} 137B for GSM8K. Among 50 randomly selected instances where the model produced the correct final answer, every generated chain of thought was both logically and mathematically sound, except for two cases where the model arrived at the correct answer by coincidence (refer to \cref{subsec:correct-chain-of-thought} and \cref{tab:appendix-gsm8k-correct-examples} for examples of accurate model-generated chains of thought). Additionally, we randomly inspected 50 samples where the model produced incorrect answers. Our analysis revealed that 46\% of these chains were nearly correct, with only minor errors such as calculator mistakes, symbol mapping errors, or a single missing reasoning step, while the remaining 54\% contained significant semantic or coherence errors (see \cref{subsec:incorrect-chain-of-thought-analysis}).

To offer some insight into why increasing model scale enhances chain-of-thought reasoning capabilities, we conducted a similar error analysis on \palm{} 62B and examined whether these errors were addressed when scaling to \palm{} 540B. The findings indicate that scaling \palm{} to 540B substantially reduces the frequency of one-step omissions and semantic understanding errors observed in the 62B model (see \cref{subsec:why-scale-helps}).

\subsection{Ablation Study}\label{subsec:math-ablation}

The demonstrated advantages of chain-of-thought prompting naturally lead to the question of whether similar performance gains can be achieved through alternative prompting methods. 
\cref{fig:bar_ablation} presents an ablation study involving three variations of chain-of-thought prompting as outlined below.

\textbf{Equation only.} One possible reason chain-of-thought prompting aids performance is that it generates the mathematical equation to be solved. To test this hypothesis, we evaluated a variant where the model is instructed to output only the mathematical equation prior to providing the final answer. 
As shown in \cref{fig:bar_ablation}, equation-only prompting yields little improvement on GSM8K, suggesting that the semantic complexity of GSM8K questions makes direct translation into an equation difficult without the intermediate natural language reasoning present in the chain of thought. Conversely, for datasets consisting of one-step or two-step problems, we observe that equation-only prompting does enhance performance, as the equation can be straightforwardly extracted from the question (see Appendix \cref{tab:ablations-arithmetic}).

\input{main_fables/ablation-bar}
\textbf{Variable compute only.} Another hypothesis is that the chain of thought enables the model to allocate more computation (i.e., intermediate tokens) to more difficult problems. To separate the impact of variable computation from chain-of-thought reasoning, we experiment with a setup where the model is prompted to output only a sequence of dots ($\ldots$) matching the number of characters in the equation required to solve the problem. This variant achieves performance similar to the baseline, indicating that variable computation alone does not account for the effectiveness of chain-of-thought prompting, and that there seems to be value in articulating intermediate steps using natural language.

\textbf{Chain of thought after answer.} Another possible advantage of chain-of-thought prompting might be that such prompts help the model better access relevant knowledge gained during pretraining. To test this, we try an alternative setting where the chain of thought prompt is provided only after the answer, isolating whether the model truly relies on the generated chain of thought to produce the final answer. This variant performs comparably to the baseline, suggesting that the sequential reasoning embodied in the chain of thought contributes benefits beyond merely triggering knowledge.

\input{main_fables/main-robustness}
\subsection{Robustness of Chain of Thought}\label{subsec:robustness}

Sensitivity to exemplars is a crucial factor for prompting methods—for example, changing the order of few-shot exemplars can cause GPT-3’s accuracy on SST-2 to vary from nearly chance level (54.3\%) to nearly state-of-the-art (93.4\%) \citep{zhao2021calibrate}. In this final subsection, we assess the robustness of chains of thought authored by different annotators. Beyond the results above, which used chains of thought by Annotator A, two additional co-authors of this paper (Annotators B and C) independently wrote chains of thought for the same few-shot exemplars (presented in \cref{sec:appendix-alternate-annotators}). Annotator A also produced another, more concise chain of thought than the original, following the style of solutions in \citet{cobbe2021training}.\footnote{For example, while the original chain of thought employs several short sentences (\textit{``There were originally 9 computers. For each of 4 days, 5 more computers were added. So 5 * 4 = 20 computers were added. 9 + 20 is 29.''}), the concise chain of thought would be \textit{``5 * 4 = 20 new computers were added. So there are 9 + 20 = 29 new computers in the server room now''}.}

\cref{fig:bar-robustness-analysis} presents the outcomes for \lamda{} 137B on the GSM8K and MAWPS datasets (additional ablation results for other datasets can be found in Appendix \cref{tab:ablations-arithmetic} / \cref{tab:ablations-commonsense-symbolic}). While there is some variation among different chain of thought annotations, which is anticipated when employing exemplar-based prompting \citep{le-scao-rush-2021-many,reynolds2021prompt,zhao2021calibrate}, every set of chain of thought prompts significantly surpasses the standard baseline. This indicates that the effectiveness of chain of thought prompting is not dependent on a specific linguistic style.

To verify that effective chain-of-thought prompting holds for other sets of exemplars, we conducted experiments using three sets of eight exemplars randomly drawn from the GSM8K training set, an independent source (the examples in this set already include reasoning steps akin to a chain of thought).\footnote{We sampled examples containing $\leq 60$ tokens to fit within our input context window and restricted the examples to $\leq 2$ steps to ensure a fair comparison with the eight exemplars we created.} \cref{fig:bar-robustness-analysis} demonstrates that these prompts performed similarly to our manually crafted exemplars, and both significantly outperformed the standard prompting approach.

Beyond robustness to different annotators, independently created chains of thought, varying exemplars, and multiple language models, we also observe that chain-of-thought prompting for arithmetic reasoning is stable with respect to different exemplar orders and varying quantities of exemplars (refer to \cref{subsec:faq-prompt-engineering}).

\section{Commonsense Reasoning}\label{sec:commonsense-reasoning}

Although chain of thought is especially effective for math word problems, its language-based character renders it suitable for a wide range of commonsense reasoning tasks, which require reasoning about physical and human interactions assuming general background knowledge. Commonsense reasoning is essential for interacting with the world and remains beyond the capabilities of current natural language understanding systems \citep{talmor2022commonsenseqa}.

\textbf{Benchmarks.}
We examine five datasets that represent a broad spectrum of commonsense reasoning challenges. The widely used \textbf{CSQA} \citep{talmor-etal-2019-commonsenseqa} presents commonsense questions about the world that involve intricate semantics and often necessitate background knowledge. \textbf{StrategyQA} \citep{geva-etal-2021-aristotle} demands that models deduce a multi-step strategy to arrive at answers. We select two targeted evaluation sets from the BIG-bench initiative \citep{bigbench}: \textbf{Date} Understanding, which requires extracting a date based on provided context, and \textbf{Sports} Understanding, which entails assessing whether a sports-related sentence is credible or not. Lastly, the \textbf{SayCan} dataset \citep{ahn2022can} involves converting a natural language instruction into a sequence of robotic actions drawn from a discrete set. \cref{fig:dataset-examples} presents examples for all datasets accompanied by chain of thought annotations.

\textbf{Prompts.}
We adopt the same experimental procedure as described in the previous section. For CSQA and StrategyQA, we randomly picked instances from the training sets and manually crafted chains of thought to serve as few-shot exemplars. Since the two BIG-bench tasks lack training sets, we selected the first ten examples from their evaluation sets as few-shot exemplars and evaluated performance on the remaining examples. For SayCan, we utilized six training examples from \citet{ahn2022can}, also manually creating chains of thought for them.

\textbf{Results.}
\cref{fig:commonsense-results} presents these findings for \palm{} (complete results for \lamda{}, GPT-3, and various model sizes are provided in \cref{tab:all-lm-commonsense}).
Across all tasks, increasing model size enhanced the effectiveness of standard prompting; applying chain-of-thought prompting yielded additional improvements, with the most significant gains observed in \palm{} 540B. Under chain-of-thought prompting, \palm{} 540B demonstrated strong performance relative to baseline models, surpassing the previous state of the art on StrategyQA (75.6\% versus 69.4\%) and exceeding the accuracy of an unaided sports enthusiast in sports understanding (95.4\% versus 84\%).
These outcomes indicate that chain-of-thought prompting can boost performance on tasks involving diverse commonsense reasoning skills (although the improvement was minimal on CSQA).

\input{main_fables/main-commonsense}

\input{main_fables/main-symbolic}
\section{Symbolic Reasoning}\label{sec:symbolic-reasoning}
Our final set of experiments addresses symbolic reasoning, which is straightforward for humans but potentially difficult for language models. We demonstrate that chain-of-thought prompting not only enables language models to tackle symbolic reasoning tasks that are hard under standard prompting, but also supports length generalization to input sequences longer than those seen in the few-shot examples.

\paragraph{Tasks.}
We employ the following two simple tasks.

\begin{itemize}[leftmargin=*,topsep=0pt]
    \itemsep0em \item \textbf{Last letter concatenation.} This task requires the model to concatenate the last letters of words in a name (for example, \textit{``Amy Brown''} $\rightarrow$ \textit{``yn''}). 
    This is a more difficult variant of first letter concatenation, a task language models can already perform without chain of thought.\footnote{We tested 10 common names using GPT-3 \texttt{davinci} and it answered all but one correctly.} We generate full names by randomly combining first and last names from the top one-thousand most common names based on name census data ({\small{\url{https://namecensus.com/}}}).
    \item \textbf{Coin flip.} This task asks the model to determine if a coin remains heads up after a series of people either flip or do not flip the coin (e.g., \textit{``A coin is heads up. Phoebe flips the coin. Osvaldo does not flip the coin. Is the coin still heads up?''} $\rightarrow$ \textit{``no''}).
\end{itemize}

Since the design of these symbolic reasoning tasks is precisely defined, for each task we evaluate on an \textit{in-domain} test set where the examples contain the same number of steps as those in the training or few-shot exemplars, as well as on an \textit{out-of-domain} (OOD) test set, where evaluation examples involve more steps than the exemplars. For the last letter concatenation task, the model is only exposed to exemplars involving names with two words, and is then tested on last letter concatenation with names containing 3 and 4 words.\footnote{For names longer than 2 words, we concatenate multiple first and last names together.} A similar approach is applied regarding the number of possible flips in the coin flip task. Our experimental setup employs the same methods and models described in the previous two sections. We once again manually craft chains of thought for the few-shot exemplars for each task, which are presented in \cref{fig:dataset-examples}.

\paragraph{Results.}
The outcomes of these in-domain and OOD evaluations are displayed in \cref{fig:symbolic-main} for \palm{}, with results for \lamda{} included in Appendix \cref{tab:all-lm-symbolic}. With \palm{} 540B, chain-of-thought prompting achieves nearly 100\% success rates (it should be noted that standard prompting already solves the coin flip task with \palm{} 540B, although not with \lamda{} 137B). It is important to recognize that these in-domain tests are considered ``toy tasks'' because the perfect solution structures are already given by the chains of thought in the few-shot exemplars; the model's job is simply to replicate the same steps with the new symbols in the test examples. Nevertheless, smaller models still fail—capability to carry out abstract manipulations on unseen symbols for these three tasks emerges only at the scale of 100B model parameters.

Regarding the OOD evaluations, standard prompting fails on both tasks. However, with chain-of-thought prompting, language models demonstrate upward scaling trends (although the performance remains lower than in the in-domain scenario). Therefore, chain-of-thought prompting enables length generalization beyond previously seen chains of thought for sufficiently large language models.

\section{Discussion}\label{sec:discussion}

We have investigated chain-of-thought prompting as a straightforward approach to evoke multi-step reasoning capabilities in large language models. Initially, we observed that chain-of-thought prompting significantly enhances performance on arithmetic reasoning, yielding gains far exceeding those of ablations and showing robustness across different annotators, exemplars, and language models (\cref{sec:arithmetic-reasoning}). Subsequently, experiments on commonsense reasoning highlighted the broad applicability of chain-of-thought reasoning due to its linguistic nature (\cref{sec:commonsense-reasoning}). Finally, we demonstrated that for symbolic reasoning tasks, chain-of-thought prompting supports OOD generalization to longer input sequences (\cref{sec:symbolic-reasoning}). Throughout all experiments, chain-of-thought reasoning was elicited simply by prompting existing off-the-shelf language models without any fine-tuning conducted during the course of this work.

The emergence of chain-of-thought reasoning as a function of model scale has been a consistent theme \citep{wei2022emergent}. For numerous reasoning tasks where standard prompting shows a flat scaling curve, chain-of-thought prompting produces markedly steeper scaling trajectories. Chain-of-thought prompting seems to broaden the range of tasks that large language models can successfully tackle—in other words, our findings highlight that standard prompting merely provides a lower bound on the capabilities of large language models. This insight likely prompts more questions than answers—for example, to what extent might reasoning ability improve with further increases in model size? What other prompting strategies could further expand the array of tasks solvable by language models?

Regarding limitations, we first note that although chain-of-thought mimics human reasoning processes, this does not establish whether the neural network is genuinely "reasoning," which remains an open question. Second, while the manual effort to augment exemplars with chains of thought is minimal in the few-shot context, such annotation costs could become prohibitive for fine-tuning (although this might be overcome through synthetic data generation or zero-shot generalization). \orange{Third, there is no assurance that the reasoning paths are always correct, which can result in both accurate and inaccurate answers; enhancing the factual reliability of language model generations is a key direction for future research \citep[][\textit{inter alia}]{rashkin2021measuring,ye2022unreliability,wiegreffe2021reframing}.} Lastly, the fact that chain-of-thought reasoning emerges only at large model scales makes deployment expensive in real-world applications; further investigations could explore methods to induce reasoning capabilities in smaller models.

\section{Related Work} This research is influenced by various fields, which we elaborate on in a more comprehensive related work section (\cref{sec:extended-related-work}). Here, we focus on two particular lines of research and the corresponding studies that are arguably most pertinent.

The first pertinent area involves leveraging intermediate steps to address reasoning challenges. \cite{ling-etal-2017-program} introduce the concept of employing natural language rationales to tackle math word problems by breaking down the solution into sequential intermediate steps. Their approach stands in notable contrast to prior work that uses formal languages for reasoning \citep{roy-2016-reasoning, chiang-chen-2019-semantically, amini-etal-2019-mathqa, chen2019neural}. Building on \citet{ling-etal-2017-program}, \citet{cobbe2021training} create a larger dataset and utilize it to fine-tune a pretrained language model rather than training from scratch. In the context of program synthesis, \cite{nye2021show} exploit language models to predict the final outputs of Python programs by initially predicting intermediate computational states line-by-line, demonstrating that this stepwise prediction method surpasses directly forecasting the final outputs.

This study also connects closely with the extensive recent literature on prompting techniques. Since the advent of few-shot prompting popularized by \citet{brown2020language}, various general strategies have enhanced model prompting capabilities, such as learning prompts automatically \citep{lester-etal-2021-power} or providing models with task instructions \citep{wei2021finetuned,sanh2021multitask,ouyang2022training}. While these methods focus on improving or supplementing the input side of prompting (e.g., instructions prepended to inputs), our approach diverges by enriching the output of language models through generating a chain of thought.

\section{Conclusions} We have investigated chain-of-thought prompting as a straightforward and widely applicable technique for boosting reasoning abilities in language models. Our experiments covering arithmetic, symbolic, and commonsense reasoning reveal that chain-of-thought reasoning emerges as a property related to model scale, enabling sufficiently large language models to solve reasoning tasks that otherwise exhibit flat performance scaling. Expanding the range of reasoning capabilities in language models may encourage further research into language-driven reasoning methods.

\end{document}